\section{Folding a layer and paying its cost}\label{sec:fold}
We work under the strong-induction hypothesis that every graph on fewer
than $n$ vertices has a packing of weight at most $C\phi(v)v$.
At an active stage $0\le i\le k$, put $w=w_i\ge w_A$ and
\[
 S=S_i,\qquad T=\sum_{j=0}^{i+2}S_j,\qquad f=\Phi(bw).
\]
Descendant pieces are grouped inside their actual DFS boxes into batches.
Write $D$ for total batch order and $A$ for the total size of the active
blocks that are contracted. The graph construction proves
\begin{equation}\label{eq:fold-masses}
 S-(N(w)+U(w))n\le A\le S+O(w)n,
       \qquad D\le(1+O(w))n.
\end{equation}
Here $N$ pays the small or unusable descendants, $U$ the discarded boundary
and grouping terms, and $O$ the overlapping support multiplicity.
Each batch has order at most $2^{bw}$. Its active blocks are pairwise
vertex-disjoint and belong to uniquely determined parent cores.

\begin{lemma}[Physical fold]\label{lem:physical}
There is an actual packing of the whole graph $G$ of weight at most
\begin{align}
 &(1+C\eta)T+\sigma S+S/2048+(B(w)+g(w))A\notag\\
 &\hspace{20mm}+Cf(D-\kappa A)+(J(w)+V(w))n.\label{eq:physical}
\end{align}
\end{lemma}
\begin{proof}
For each batch, retain one original representative of every quotient edge;
classify the original edges lost by contraction as internal pairs or
matching requests. This produces a simple quotient. An active block saves
at least $\kappa$ times its order in the quotient order, so the summed
quotient orders are at most $D-\kappa A$. Each quotient has order at most
$2^{bw}\le n^b<n$, since $n>1$ and $b<1$. Apply the induction hypothesis
and monotonicity of $\phi$ to obtain quotient packing weight at most
$Cf(D-\kappa A)$.

Use the joint routing construction of Section~\ref{sec:routing} on each
selected core. The single original-edge-endpoint injection covers low
pairs, paired matching edges, and switches in quotient cycles simultaneously.
The resulting four families are: completed core graphs, closed low groups,
closed matching groups, and lifted quotient packings. Uncovered quotient
edges keep their chosen original singleton representatives. Unused routing
paths are not deleted from the core.

The completed core and low-group families have combined weight at most
\[
 (1+\sigma+C\eta+1/2048)S
                  +\text{(prepared pair count)}/Q(w).
\]
This bound uses the high-centre credit of Lemma~\ref{lem:completion}
before any mass upper bound is substituted. The matching groups and the
prepared-pair fee together are at most $(B(w)+g(w))A$, and the soft routing
charge for lifted quotient cycles is at most $V(w)n$. The quotient packing
contributes the term already obtained above.

The four families are pairwise edge-disjoint and their union is exactly
the selected cores and the full descendant frontier. To see the importance
of both assertions, separate original non-routing edges from the paths
deleted during routing. Every original edge has exactly one of the four
owners by the quotient classification. Every deleted path has exactly its
request owner; joint routing makes different owners edge-disjoint. An edge
cannot belong to both classes because the core remainders were defined by
deleting precisely the selected paths. This is the partition proved in
\mcode{ActualFoldDisjoint12} and \mcode{ActualFoldCover12}.

Finally insert the other prefix cores, of total mass $T-S$, at cost
$(1+C\eta)(T-S)$, and the prefix long cycles at cost at most $J(w)n$.
The prefix partition from Section~\ref{sec:layers} now covers the whole
original graph. Adding the displayed costs proves~\eqref{eq:physical}.
\end{proof}

Define the coefficient of the selected core mass by
\begin{equation}\label{eq:fold-coefficient}
 \Gamma(P,w)=(1+C\eta)P+\sigma+B(w)+g(w)+\frac1{2048}
                                     -C\kappa\Phi(bw).
\end{equation}
\begin{lemma}[Fold ledger]\label{lem:ledger}
If $T\le PS$, then the right-hand side of~\eqref{eq:physical} is at most
\begin{equation}\label{eq:paid-fold}
                Cf n+\Gamma(P,w)S+CH(w)n.
\end{equation}
\end{lemma}
\begin{proof}
Write the terms involving $A$ as $BA-(C\kappa f-g)A$.
The second coefficient is nonnegative by~\eqref{eq:fee-endpoints}.
Consequently the \emph{lower} bound on $A$ in~\eqref{eq:fold-masses}
is used in its negative term, and the upper bound is used only in $BA$.
The resulting error coefficient of $n$ is
\[
 CfO+(C\kappa f-g)(N+U)+BO+J+V.
\]
All losses are nonnegative, $0<f\le1$, $g\ge0$, and $B\le B_0$, so this
is at most $CO+C\kappa(N+U)+B_0O+J+V=CH$.
The coefficient of $S$ is exactly~\eqref{eq:fold-coefficient}.
This proves~\eqref{eq:paid-fold}. In particular, one does not introduce
an extra $gO$ term by applying an upper bound to the negative group term.
\end{proof}

For each of the $40$ phase intervals $[a,d]$, the certificates supply
positive thresholds $P_t$ for all $t\ge3$. They prove
\begin{equation}\label{eq:fold-phase-payment}
 \Gamma(P_t,w)\le0\quad\text{whenever}\quad
                   w\ge a\theta_2^{-(t-3)}.
\end{equation}
There are explicit thresholds for $3\le t\le16$ and a constant threshold
for $t\ge16$. The scalar proof uses the monotonicity of $B,g,\Phi$;
the constant tail therefore requires only its first scale endpoint.

If $i$ is a stage with $S_i>0$ and
\begin{equation}\label{eq:fold-test}
                   \sum_{j=0}^{i+2}S_j\le P_{m_0-i}S_i,
\end{equation}
then $t=m_0-i$ and~\eqref{eq:schedule} give
$w_i=W\theta_2^{-(t-3)}\ge a\theta_2^{-(t-3)}$.
Lemmas~\ref{lem:physical} and~\ref{lem:ledger}, followed by
\eqref{eq:fold-phase-payment} and~\eqref{eq:potential-payment}, give
\[
        \wt_\sigma(\mathcal P)\le Cf n+CH(w_i)n\le C\phi(n)n.
\]
This completes every branch in which an active fold exists.
